\section{Transport along actual arcs}\label{sec:actual-arcs}

We first transfer the fixed-connector rule to arcs in the original random
configuration, retaining their assigned prepared endpoint states. We
then pass from arc contributions to smooth weights. The second step
requires a separate control of signed sums and completes bulk transport.

Fix a countable sequence $h\downarrow0$.  Starting with any microscopic
sequence, make the subsequence choices of Section~\ref{sec:affine}
simultaneously for these $h$'s.  The grid translates of the menu, their
representatives, and their integer winding classes form countable lists;
admissions can be stabilized for each $h$ because there are finitely many
types.  All microscopic limits below are taken on this one further
subsequence.  In particular the scalars $a_h^Q$ are bounded uniformly in
$h$ and $a_h^V=1$.  Comparisons involving an empty set of admitted types
are vacuous.

We continue to work in a bulk patch containing $\overline B(x,6R)$.
Every qualified pin is used with a specified fitting test, type, ranks,
and the prepared state assigned in Section~\ref{sec:pins}.  No
comparison below changes that assigned state.

For two ordered pins $s,t$ on $P$, write
$D_s=B_s=\overline B(z_s,m_sh)$ for the previously defined placed
cutoff ball at $s$, and similarly $D_t=B_t$ at $t$. Define
\[
 I_X(P;s,t)
   =X_P(\ind_{D_s})
     +\sum_{\substack{e\text{ on the forward arc }p_s\to p_t\\
               \text{right tag of }e\notin D_s\cup D_t}}
                   X_P(e),\qquad X\in\{Q,V\}.
\]
Here $X_P(e)$ denotes the contribution of the oriented edge with its
right tagged triangle.  The first term is the \emph{entire} start
cutoff mass, not only the portion after $p_s$.  This convention will
make concatenation exact.

\subsection{Two conditional-measure and locality facts}

\begin{lemma}[A dominated forcing kernel]\label{lemma:dominated-kernel}
Let $\nu$ be a probability measure on configurations, let $Z$ be
retained data, and let $\nu_z$ be a regular conditional law given $Z=z$.
Suppose $H$ is a measurable set of retained data and
$p(z)=\nu_z(W)\ge q>0$ on $H$.  Let $K$ fix $Z$ and satisfy
$\nu_zK_z=\nu_z(\cdot\mid W)$ there.  If an unnormalized
subprobability $\lambda$ satisfies $\lambda\le C\nu$ and is supported
over $H$, then
\[
                         \lambda K\le (C/q)\nu.
\]
For the additional monotone assertion, suppose that, conditional on
$Z=z$, the remaining coordinates have law
$\nu_z=\pi_z\otimes\tau_z$, where $\pi_z$ is a product spin law on
a finite, $Z$-measurable set $I(z)$ of exterior sites and $\tau_z$
is the law of the complementary data. Suppose also that $W$ depends
only on $Z$ and the spins in $I(Z)$ and is increasing in white on
those spins. Then $K$ can be chosen to recolor only existing black
sites in $I(Z)$ white and to fix $Z$ and every complementary coordinate.
\end{lemma}

\begin{proof}
For every nonnegative measurable function $g$, positivity of a kernel
and disintegration give
\begin{align*}
 \int\lambda(\dd\omega)Kg(\omega)
 &\le C\int\ind_H(Z(\omega))\nu(\dd\omega)Kg(\omega)\\
 &=C\int_H\nu^Z(\dd z)\,
       \frac{1}{p(z)}\int\ind_W(\omega')g(\omega')
                                      \nu_z(\dd\omega')\\
 &\le(C/q)\int g\,\dd\nu.
\end{align*}
No assumption is made about the conditional law of $\lambda$ given
$Z$. Under the additional hypothesis, association of $\pi_z$ shows
that $\pi_z(\cdot\mid W)$ stochastically dominates $\pi_z$ in the
white order. The finite ordered coupling construction in
Section~\ref{sec:pins}, applied on each retained-data fiber, gives
a measurable white-adding kernel for this block. Independence of
the complementary data lets us use the identity there: the output
law is $\pi_z(\cdot\mid W)\otimes\tau_z=\nu_z(\cdot\mid W)$.
\end{proof}

\par\ifdim\dimexpr\pagegoal-\pagetotal\relax<6\baselineskip\newpage\fi
\begin{lemma}[Retention with collars]\label{lemma:retained-collar}
Fix $h$ and finitely many deterministic path envelopes and state disks.
Suppose the point data are kept on a deterministic Borel set containing
an open collar of every envelope and of the state portions under
comparison.  Set $\rho_\eps=\sqrt{h\eps}$.  For small $\eps$, every
Delaunay triangle of circumradius at most $\rho_\eps$ meeting these
envelopes, together with its incident edge and tag data there, is
determined by the kept point restriction.  Thus two configurations
with the same restriction have the same such local triangulations.

In a fixed two-plug comparison one can keep the two endpoint states and
the whole exterior sector of the used connector.  The inset paths of
the cycle experiment and any sufficiently closely tracking actual
arc, including their inner second-gate paths, then have the common
triangulations just described on the ordinary mesh-good events.  The
complementary connector can be resampled without altering them.
\end{lemma}

\begin{proof}
A Delaunay triangle is determined by its three vertices and emptiness
of its circumdisk, apart from the null degeneracies excluded in our
model.  If its circumradius is at most $\rho_\eps$ and it meets a fixed
envelope, its circumdisk lies within a constant multiple of
$\rho_\eps$ of that envelope.  This is eventually inside its kept
collar.  The same observation determines all retained triangles
incident to its local edges.  On mesh-good events all triangles
needed for the ordinary local paths and their right tags satisfy this
bound, by Lemma~\ref{lemma:mesh}.  This proves the first assertion
without any assumption on the area of the boundaries of the kept
sets.

For the second assertion, the paths in the cycle experiment lie in
fixed inset layers of the embedded annular band.  Outside states these
layers, the tube-test footprints, and the joining seams have positive
margins to the lateral boundaries of the used sector.  Inside states,
the relevant second-gate paths and cutoff triangles lie in a compact
subset of $O$, since the gate annuli and cutoff balls have the
prescribed radial margins.  The same envelopes contain the actual
paths when the track fit is sufficiently close.  Their exterior
portions and state portions therefore have collars within the kept
union, as also in Lemma~\ref{lemma:fixed-collar}.  The other sector
queries disjoint exterior restrictions.  Deterministic Borel
allocation of boundaries suffices; no smooth-boundary convention is
needed.  The first assertion now applies in both configurations.
\end{proof}

\subsection{Winding and a finite deterministic family of comparisons}

\begin{lemma}[Winding bound]\label{lemma:winding}
Except on events of probability tending to zero in the ordered limit,
every subarc of a rim in $B(x,3R)$ that stays at distance at least $h$
from a grid center in use has absolute net change of continuous
argument about that center at most
$C_R(1+|\log h|^2)$.
\end{lemma}

\begin{proof}
Fix a ray from the center and divide its relevant distances into
dyadic bins $d\asymp 2^jh$, up to scale $O(R)$.  In each bin a bounded
number of balls of radius a small fixed fraction of $d$ cover possible
ray passages.  Give each ball a buffered outer ball of comparable
radius, all within the bulk patch.  A passage off the ray followed
by a quarter turn leaves one of these outer balls, and so gives a
monochromatic crossing of its fixed-ratio annulus.  Successive full
turns contributing to a net argument change can be selected by first
hits of successive levels of a lifted argument.  Their portions
between the relevant passage and departure times are disjoint on the
simple rim; trimming with the fixed margins yields vertex-disjoint
annular crossings on fine mesh.

On crossing-good geometry an opposite-color circuit in such a test
annulus has probability at least a fixed $q>0$.  Hence one
monochromatic crossing has probability at most $1-q$, and conditional
BK bounds the probability of $k$ disjoint crossings by $(1-q)^k$.
There are $O_R(h^{-2})$ centers and $O(1+|\log h|)$ bins per center.
At each fixed $h$ the geometry exceptions for this finite list vanish
as $\eps\to0$.  A union bound with
$k=C_R(1+|\log h|)$ leaves probability tending to zero as $h\to0$.
Summing over the bins bounds the net number of turns by
$C_R(1+|\log h|^2)$, and the incomplete initial and final turns add
only a constant.
\end{proof}

Fix $h$ for the moment and an arbitrarily small probability budget
$\zeta>0$.  Impose the local inverse-diameter tests from
Section~\ref{sec:pins} throughout $\overline B(x,3R)$, to a finite
number of accuracy levels to be chosen below, at total limiting cost
less than $\zeta$.  Track a qualified arc from $s$ to $t$ together
with its extension from the start of $K_s$ to the end of $K_t$, saving
both marked endpoint profiles.  Assume the original arc lies in
$\overline B(x,2R)$, $|z_s-z_t|\ge R/3$, and its complementary rim
reaches outside $B(x,4R)$.  For small $h$ the endpoint neighborhoods
of size $10Lh$ are disjoint and occur in the required order, so the
extended arc lies in $B(x,5R/2)$ and its complement still makes a
fixed large departure.

The proof of Lemma~\ref{lemma:profile-compactness}, in this fixed
physical region, applies to these extended marked arcs.  Along
sequences satisfying increasingly many constraints it gives simple
Fr\'echet limits, simultaneously with their endpoint profile limits.
At fixed $h$ the placements, templates, and retained ranks lie in
finite lists.  Pass to fixed choices.  The endpoint limits are in the
closures of their assigned fit neighborhoods.  A visit near an end
obstacle not belonging to its prescribed $K$ is included in that
profile's $E$: the relevant neighborhood is well inside the radius
$50Lh$ truncation.

Lemma~\ref{lemma:exact-fit} supplies independently supported permitted
changes of the two end charts, making the limiting extended arc the
exact centerline through both changed obstacles and for a further
collar on either side.  It avoids the obstacles elsewhere.  The
band-and-collar construction of Section~\ref{sec:plugs} now fits a
connector to the intervening simple arc, with the full prescribed
port sections, inside the bulk patch.  A complementary connector
exists in the plane; it need not lie in the patch.  The two-twist
classification applies: a net argument along a limiting arc separated
from a center persists under close Fr\'echet approximation.  Thus
Lemma~\ref{lemma:winding} bounds the two winding integers by
$C_R(1+|\log h|^2)$, up to a bounded contribution from the fixed stubs
and allowed chart changes.  Since $|z_s-z_t|/h\ge R/(3h)$, the
distance requirement in Proposition~\ref{prop:endpoint-potentials}
holds for small $h$.

This fixed connector works on an open neighborhood of the limiting
track data.  In particular the actual root vertices near nominal
$u_s=1.7$ and $u_t=-1.7$ are joined by the actual rim within its
black tube.  The corresponding compact portion of the limiting
centerline is interior to that tube and is well beyond the
longitudinal cutoffs $u_s=0.5$ and $u_t=-0.5$.  Small allowed chart
changes preserve the nominal root positions and windows.  The
product-band seams are interior, so a sufficiently close ordered fit
and fine mesh preserve all these margins, including retention of the
required paths.

There is a further gate-order issue.  Use the first and second
circuits selected from the prepared states.  Every contact with the
actual rim lies on its isolated central traversal and remains black
after preparation.  Each selected circuit has at most two black
vertices in that prepared state, so it has at most two contacts with
the original rim.  Passage from inside to outside requires two;
there are therefore exactly two contacts, traversed once.  This
argument applies even if a selected circuit differs from the
original gate choice.  A full chart slab about
zero lies in the inner gate ball.  The limiting track enters this
slab from the negative side and leaves on the positive side; its
contacts with a surrounding gate lie away from the slab.  On a
sufficiently close ordered fit their signs cannot change.  Thus the
first exit at $s$ and the first entry at $t$ are the designated
positive and negative gates.  The intervening path remains in the
used sector, with room in its interior together with the endpoint
states, and avoids the opposing layer.  It cannot recross a central
section while outside the corresponding first gate.  The inner
second-gate paths also track the local centerline inside their
states.  This argument uses the progress through a slab, not a
transversal crossing of a discrete edge with a curve, and is uniform
over the possible locations of the circuits in their annuli.

Compactness and the finite-list argument in
Lemma~\ref{lemma:exact-fit} now give a finite family of these
deterministic patterns whose open neighborhoods cover all tracks
satisfying a sufficiently long finite constraint list and the winding
bound.  Patterns are chosen from limiting deterministic track data
before the random comparison; a new band is not selected after
observing the colors.  Only finitely many placed patterns, endpoint
types, and rank choices occur at this fixed $h$.  The open track
tests and existential path conditions are measurable: graph paths
are countable given the locally finite graph, and the marked fits
use the usual Fr\'echet open sets.

\subsection{The comparison with a cycle experiment}

\begin{proposition}[Transport on a real arc]\label{prop:real-arc}
Simultaneously for $X=Q,V$ and every qualified ordered arc from $s$
to $t$ contained in $\overline B(x,2R)$, with
$|z_s-z_t|\ge R/3$ and complementary rim reaching outside
$B(x,4R)$, one has
\begin{equation}\label{actual:arc}
 \left|I_X(P;s,t)-a_h^X(z_t-z_s)
       -\phi_s^{X,+}(\sigma'_s)-\phi_t^{X,-}(\sigma'_t)\right|
       \le C_Rh(1+|\log h|^2)
\end{equation}
outside events whose probability tends to zero in the ordered limit.
The constant is independent of the complexity and conditioning
constants of the individual fixed comparison patterns.
\end{proposition}

\begin{proof}
We construct one comparison for all qualifying arcs of each pattern
in the finite deterministic cover.  Write $\xi$ for the original input,
with law $\mathsf P$, including the bulk/full-plane coupling and all
assigned local preparation seeds.  Let $E_0$ be the common input event
on which patch agreement and the covering and winding conclusions hold,
intersected with
the original mesh and continuation-geometry conditions of all patterns.
These additional conditions have total failure probability $o(1)$ at
fixed $h$.  The covering exception will be paid once, through
$\mathsf P(E_0^c)$.

Fix a placed connector $e:s\to t$, a complementary connector $c$, and
their endpoint choices.  Let $\mathscr C_e(\xi)$ be the collection of
all qualifying ordered rim arcs of this pattern, including their assigned
preparation outcomes.  Put
\[
 A_e=\{\mathscr C_e(\xi)\ne\varnothing\},\qquad A_{e,0}=A_e\cap E_0.
\]
We retain the original input throughout and work with the
\emph{unnormalized} restriction $\mathsf P|_{A_{e,0}}$; its mass is
never divided by the probability of a fit.  We describe the construction
for black arcs; for white arcs exchange the colors.

\paragraph{Prepare the endpoints and force the white connection.}
Let $F_e(\xi)$ be the underlying plane configuration after the two
specified local white-adding preparations, with their already assigned
seeds.  Their supports are disjoint.  Before restriction, independence
of the seeds conditional on geometry and the local preparation bounds
give a configuration marginal bounded by $C_{\rm prep}\nu$.  Hence
\[
 \lambda:=F_{e*}(\mathsf P|_{A_{e,0}})\le C_{\rm prep}\nu.
\]
Restriction only deletes mass, even when it inspects actual arcs,
other preparations, seeds, and patch data.  The endpoint states of
$F_e(\xi)$ are precisely the assigned $\sigma'_s(\xi),\sigma'_t(\xi)$.

Define $Z$ on the ordinary configuration space to be the coordinate
map consisting of the full point geometry $\eta$ and the colored
configuration restrictions to the two state regions $O_s,O_t$.
On a prepared output these last two coordinates have the assigned
values $\sigma'_s,\sigma'_t$. Thus $\nu(\cdot\mid Z)$ refers to
ordinary coordinate restrictions, not to conditioning on the outcome
of a preparation algorithm under its input law. Let $W$ be the white
part of $S_e$. Under the ordinary law conditional on $Z$, the exterior
colors are independent fair bits. Each needed white side event has probability at least
$\beta$ by qualification.  These two endpoint events use disjoint
exterior bits.  The white continuation network uses the buffered
exterior corridors of Lemma~\ref{lemma:bridge}.  On its good point
geometry, association in the white order therefore gives
\begin{equation}\label{actual:quenched-force}
                  \nu(W\mid Z)\ge q_{\rm cont}\beta^2=:q_0>0.
\end{equation}
The local qualification is unchanged by additionally conditioning on
the rest of $\eta$, since its retained side tests query only their
specified local point restrictions.  The continuation-geometry conditions
were included in $E_0$ and have vanishing failure probability; the
preparations keep the point geometry fixed.

Let $H$ be the resulting set of $Z$ values on which
\eqref{actual:quenched-force} holds. For this fixed pattern let
$I=I(\eta)$ contain \emph{all potential exterior sites} in its
buffered white footprint on which $W$ can depend. This finite block
is selected from the geometry before colors or witnesses are inspected;
it is not an adaptive query transcript. Every site in it is separated
from the permitted black track layer. Conditional on $Z$, its fair
product spin law is independent of all complementary exterior spins,
and $W$ depends only on $Z$ and this block. Equivalently, with
$\widetilde Z=(Z,\sigma_{\rm ext\setminus I})$, these spins remain
independent fair bits and
$\nu(W\mid\widetilde Z)=\nu(W\mid Z)\ge q_0$ over $H$.
The monotone clause of Lemma~\ref{lemma:dominated-kernel} therefore
applies to $\lambda$, which is supported over $H$. Use its kernel
$K_e^W$ to add white only in $I$, fixing every other
coordinate. The lemma yields
\[
                 \lambda K_e^W\le(C_{\rm prep}/q_0)\nu.
\]
This conclusion does not assume that the selected input has ordinary
conditional exterior spins.  Its conditional law may be arbitrary;
the subprobability domination is what permits the forcing.  The
kernel only removes black sites in the separated white footprint.
It preserves every track in $\mathscr C_e(\xi)$, with its buffered
collars and fixed states.  Thus the output satisfies $S_e$.

\paragraph{Complete the configuration while retaining the used sector.}
Keep the endpoint states and the whole exterior $e$-sector datum,
denoted together by $M_e$.  Under the ordinary kept-data law
$\nu_{M_e}$, the corresponding marginal $\rho_{M_e}$ of the
two-plug cycle experiment has density
\begin{equation}\label{actual:cycle-density}
 \frac{\dd\rho_{M_e}}{\dd\nu_{M_e}}
 =\frac{\ind_{G_s}\ind_{G_t}\ind_{S_e}}
 {\nu(G_s)\nu(G_t)\nu(S_e\mid\sigma'_s,\sigma'_t)}.
\end{equation}
The denominator is positive by admission and \eqref{plug:success},
and is at most one.  The post-forcing marginal is supported on the
three events in the numerator and dominated by
$(C_{\rm prep}/q_0)\nu_{M_e}$; it is therefore dominated by
$(C_{\rm prep}/q_0)\rho_{M_e}$.  Apply to both
measures the same completion kernel: sample the other exterior
connector region, geometry and colors, conditional on $S_c$ and the
two states, and sample the remaining plane ordinarily.  This keeps
$M_e$ and the collars of the compared paths unchanged.  Denote the
composition of $K_e^W$ with this completion kernel by $\mathcal K_e$,
extending it arbitrarily off the relevant support.  On the original
input and final configuration together, define
\[
 \Pi_e(\dd\xi,\dd\omega')
   =\ind_{A_{e,0}}(\xi)\mathsf P(\dd\xi)
             \mathcal K_e(F_e(\xi),\dd\omega').
\]
Its input marginal is exactly $\mathsf P|_{A_{e,0}}$.  The construction
gives the two properties needed below:
\[
 \begin{aligned}
 &(\operatorname{pr}_{\rm out})_*\Pi_e\le D_e\rho_e,\\
 &(\sigma_s(\omega'),\sigma_t(\omega'))
       =(\sigma'_s(\xi),\sigma'_t(\xi))
       \quad\Pi_e\text{-almost surely}.
 \end{aligned}
\]
Here $\rho_e$ is the two-plug cycle law and $D_e<\infty$ is fixed
at fixed $h$ and pattern.  The first follows by applying the same
completion kernel to the dominated kept-data marginals; the second
holds because both kernels fix the endpoint states.  No product
domination of the joint law is required.

\paragraph{Compare every fitting arc to one output rim.}
Choose one experiment rim $P'_e=P'_e(\omega')$ by a measurable
rule depending only on the output.  There is one event
$\mathcal G_e$ in this joint space, independent of a choice of
$\gamma\in\mathscr C_e(\xi)$, such that
$\Pi_e(\mathcal G_e^c)=o(1)$ and the following path comparisons
hold.  Use the mesh and shielding conditions on the deterministic
state and sector envelopes in both configurations, and the
mesh conditions for the confined cycle-rim construction.
Their original failure mass is bounded by the original mesh
probability; their final failure mass is bounded by $D_e$ times
the cycle mesh probability.  The strict track and gate-order
margins hold for every member of $\mathscr C_e(\xi)$ by the one
fixed pattern's fit test.  Thus the good event and the output
kernel are not chosen separately for different arcs.

On $\mathcal G_e$, the first-exit to first-entry path of $P'_e$
and its inner second-gate paths are black in the original graph.
Colors only changed from black to white before completion,
and completion resamples outside their kept envelopes.
Conversely, the corresponding portions of every input arc in
$\mathscr C_e(\xi)$ remain black in the final graph: the
preparations and white forcing avoid their permitted track layer,
and completion keeps the whole used sector and endpoint states.
Lemma~\ref{lemma:retained-collar} gives the common complete
incident triangulations and gate fans. Since $b_2<a_3$, the first
and second gates lie strictly inside the unchanged third circuit.
Third-gate isolation puts every original-rim contact with a gate selected from the
prepared state on the untouched central traversal.  The selected
gate has at most two path-color vertices, while each fitting rim
visits both sides, so every such rim uses precisely those two
contacts.  The output rim does too, and ordered passage through
the central slab gives the same signs; this also covers a selected
gate differing from the originally exhibited gate.
Lemma~\ref{lemma:gate-consistency} now
identifies the inner second-gate paths and the outside path,
using the inside and outside starting fans, respectively.
At a first divergence a mutually black compared path would give
a forbidden bypass up to its first subsequent contact with the
whole other rim.  This argument applies in each graph separately
and to every input arc with the same $P'_e(\omega')$; the original
rim need not remain a complete rim after modification.

Set
\[
 H_e^X(\omega')
  =\widehat m_s(P'_e(\omega'))+\widehat j_e(P'_e(\omega')),
                           \qquad X=Q,V.
\]
All cutoff tags belong to the common inner second-gate paths.
The marks and inner first-gate paths lie inside the cutoffs, and
the outside-cutoff mass lies on the common outside path.  Thus
on the same event $\mathcal G_e$,
\begin{equation}\label{actual:exact-comparison}
 I_X(\gamma;\xi)=H_e^X(\omega')
 \quad\text{for all }\gamma\in\mathscr C_e(\xi),\qquad X=Q,V,
\end{equation}
where $I_X(\gamma;\xi)$ is the original $I_X(P;s,t)$.
In particular the right side is one output random variable,
not a separately coupled variable for each input arc.

\paragraph{Return the estimate to the original input.}
For endpoint states $\sigma=(\sigma_s,\sigma_t)$ write
\[
 \Theta_e^X(\sigma)=a_h^X(z_t-z_s)
        +\phi_s^{X,+}(\sigma_s)+\phi_t^{X,-}(\sigma_t).
\]
Separate the cycle-mass error from the retained-state remainder:
\[
 \begin{aligned}
 C_e^\eps(\omega')
   &=\max_X|H_e^X(\omega')
               -i_e^\eps(\sigma_s(\omega'),\sigma_t(\omega'))|,\\
 R_e^\eps(\sigma)
   &=\max_X|i_e^\eps(\sigma)-\Theta_e^X(\sigma)-b_e^{X,\eps}|.
 \end{aligned}
\]
Here $i_e^\eps$ in each term is the function for field $X$, and
$b_e^{X,\eps}$ is its deterministic bias.
Equations~\eqref{plug:cycle} and \eqref{affine:transport} give
\[
 \EE_{\rho_e}C_e^\eps=o(1),\qquad
 \EE_{\mu_s\otimes\mu_t}R_e^\eps=o(1).
\]
The input law of its two prepared states, restricted to
$A_{e,0}$, is bounded by $C_{\rm prep}\mu_s\otimes\mu_t$.
Indeed the endpoint marginal of $\lambda$ is bounded by
$C_{\rm prep}$ times the independent ordinary state product and is supported on
$G_s\cap G_t$; restricting that product gives
$\nu(G_s)\nu(G_t)\mu_s\otimes\mu_t\le\mu_s\otimes\mu_t$.
The pathwise retained-state equality makes
$\Theta_e^X(\sigma'_s(\xi),\sigma'_t(\xi))$ exactly the target
computed from the output states.

Let $B_e^\eps=\max_X|b_e^{X,\eps}|$ and fix $u>0$.
On $\mathcal G_e$ the error of every fitting input arc is at most
\[
 B_e^\eps+C_e^\eps(\omega')
             +R_e^\eps(\sigma'_s(\xi),\sigma'_t(\xi)).
\]
Consequently the original-input probability, restricted to
$A_{e,0}$, that any such arc has error greater than
$B_e^\eps+2u$ is at most
\[
 \Pi_e(\mathcal G_e^c)
 +D_e\rho_e(C_e^\eps>u)
 +C_{\rm prep}(\mu_s\otimes\mu_t)(R_e^\eps>u)=o(1).
\]
This is a projection of a pointwise assertion on one joint space;
no union over input arcs is taken.

On the common event $E_0$, every eligible arc belongs to at least
one of the finitely many patterns.  Apply the preceding estimate
to each pattern on the same original input space and take a
finite union.  The total input failure probability is bounded by
$\mathsf P(E_0^c)$ plus the sum of the three displayed error terms
over the fixed $N_h$ patterns.  The common covering exception is
therefore paid once; no term counts individual arcs.  All pattern
constants and $N_h$ are fixed before the microscopic limit.
Finally $\limsup_{\eps\to0}B_e^\eps\le Ch(1+k)$ and
$k\le C_R(1+|\log h|^2)$.  Taking $u=h$ absorbs the residual
tolerance into \eqref{actual:arc}.

At fixed $h$, the lower constants $q_0$, the cycle domination, and
the size of the cover can be arbitrarily small or large positive
constants.  They multiply errors that vanish as $\eps\to0$ at this
fixed $h$.  The deterministic $Ch(1+k)$ bias and the $Ch$ state
norms, in contrast, come from the fixed normalized short templates
of Section~\ref{sec:affine} and have uniform constants.  Since the
extra finite-cover budget $\zeta$ is arbitrary, it can be chosen
tending to zero with $h$.  The winding exceptions also tend to zero.
This proves \eqref{actual:arc} in the stated order of limits.
\end{proof}

\subsection{Additive chunks and the endpoint-potential sum}

Let $f$ be a fixed smooth function supported in $B(x,R)$.  For each
rim $P$ reaching outside $B(x,4R)$, consider its maximal excursion
intervals in $B(x,2R)$ that meet $B(x,3R/2)$.  On the density event
choose pins $s_0,\ldots,s_l$ in order in each interval so that the
initial and final unmarked portions have diameter $O(r_h)$ and
every consecutive marked subarc has diameter between $r_h$ and
$Cr_h$.  Here is an explicit choice.  Take a pin on an initial
portion reaching displacement $2r_h$.  From a chosen pin, if the
remaining excursion reaches displacement $4r_h$, apply
Proposition~\ref{prop:dense-pins} to the portion between its first
hits of displacement $2r_h$ and $4r_h$ and choose the next pin there.
Otherwise stop.  Density applies to that portion because its
diameter is at least $2r_h$.  Closure of $B(x,2R)$ may be used for
the endpoints.  The construction gives all the asserted diameter
bounds.  Set $F_i=f(z_{s_i})$.  Then $F_0=F_l=0$,
$|F_{i+1}-F_i|\le C_f r_h$, and a nonzero change places both pins
near the support, a fixed macroscopic distance from $s_0$.

For distinct chosen pins within one interval, both directions along
the rim include a departure of diameter at least $r_h$, or a
departure outside the bulk disk.  Thus neither mark belongs to the
other's third-gate arc $\Delta$, whose diameter is at most $2Lh$.
Their cutoff balls are disjoint.  Indeed if two overlapped, choose
the plug with larger $b_2$.  Since $m<b_2/10$ for each type and
$a_3>100b_2$, the other mark would lie strictly inside this plug's
third circuit and hence on its short trapped arc, a contradiction.

At an intermediate pin all its cutoff-tagged edges occur between
the preceding and following pins: their vertices belong, on fine
mesh, to the trapped $\Delta$ containing that pin, and this arc
avoids both neighbors.  At either end the corresponding statement
holds on the one neighboring side.  The first-gate paths lie inside
the cutoffs, so the convention defining $I_X$ gives the exact identity
\begin{equation}\label{actual:concatenation}
 \sum_{i=0}^{j-1}I_X(P;s_i,s_{i+1})=I_X(P;s_0,s_j),
                      \qquad 1\le j\le l.
\end{equation}
In particular, the full start cutoff and the omitted terminal cutoff
are essential parts of this equality.

Assign coefficient $F_i$ to the chunk $I_X(P;s_i,s_{i+1})$, and
coefficient zero to all other tags; denote the resulting weighted
sum by $\mathcal A_X(P)$.  Near the support every central isolating
arc lies in the same excursion, so full-cutoff reassignment cannot
mix two excursion intervals.  At end collars the coefficients are
zero and overlaps do not matter.  Every support tag is covered,
and its assigned coefficient differs from its value of $f$ by
$O_f(r_h)$, using the chunk diameter bound and fine mesh.

At active coefficient changes there is no repeated grid location
on one rim, even between different excursion intervals.  A repeat
would put the two marks inside the third circuit at that location.
Within an interval the preceding two-departure argument applies;
between intervals both directions reach the outer excursion collar.
Either possibility contradicts the single short trapped arc.
Allowing finitely many types at a location changes ensuing counts
only by a menu-dependent constant.

Summation by parts using \eqref{actual:concatenation} gives, on each
excursion,
\[
 \sum_{i=0}^{l-1}F_iI_X(P;s_i,s_{i+1})
       =\sum_{j=1}^{l}(F_{j-1}-F_j)I_X(P;s_0,s_j).
\]
For active changes the hypotheses of Proposition~\ref{prop:real-arc}
hold.  Subtract $a_h^Q$ times the identity for $V$ from that for
$Q$.  The displacement terms cancel since $a_h^V=1$, and the base
potentials cancel because
$\sum_j(F_{j-1}-F_j)=F_0-F_l=0$.  Only terminal potentials and
the bounded comparison errors remain.

We record the aggregate bounds in the original probability space,
with the jointly assigned local preparations. Let $N$ count all grid
locations near the support whose fixed $Ch$-neighborhood meets any
macroscopically reaching rim; set it to zero on the geometry exception
allowed for a bulk patch. This count dominates all possible candidate
mark locations, independently of the later pin selection or fitting
events. Each counted visit gives a one-arm event from scale $h$ to
a fixed macroscopic distance. Thus
\[
                 \limsup_{\eps\to0}\EE N\le C_R h^{-2+c}.
\]
For each location and admitted type define
$\psi^- =\phi^{Q,-}-a_h^Q\phi^{V,-}$ on that type's usable-state
event $G_s$, and zero outside $G_s$. Admission means the deterministic
inequality $\nu_s(G_s)\ge\beta'$. Its prepared-state square expectation is at most $Ch^2$:
the individual preparation law is dominated by $C\nu$, and the
$L^2(\mu)$ norm of each centered potential is at most $Ch$.
Consequently
\[
 T_h:=\sum_{\substack{\text{locations and types}\\
                         \text{in the fixed region}}}
              |\psi^-(\sigma')|^2,
       \qquad \limsup_{\eps\to0}\EE T_h\le C_R.
\]
This uses no independence of different candidate preparations.
For every adaptive pin choice on a single rim the sum of squared
selected potentials is bounded by $T_h$, by non-repetition and the
finite type list.  Cauchy--Schwarz therefore controls the maximum
over all rims by the same random bound.  Off the preceding
probability exceptions, the limiting upper expectation of the
aggregate error is at most
\begin{equation}\label{actual:coefficient-cost}
 C_{f,R}\left[
 r_hh^{-1+c}(1+|\log h|^2)+r_hh^{-1+c/2}\right].
\end{equation}
The first term is $C r_hh(1+|\log h|^2)\EE N$; the second follows
from $Cr_h\EE\sqrt{NT_h}$.  The long-connector forcing densities
from the proof of Proposition~\ref{prop:real-arc} do not enter this
estimate: the potentials here are evaluated in the original assigned
local preparations, whose uniform density constants were fixed in
Section~\ref{sec:pins}.

\subsection{Signed fields and mixed atoms}

Pointwise closeness of assigned coefficients to $f$ does not itself
control a signed field.  We now prove the needed replacement.

\begin{lemma}[Mixed-atom replacement]\label{lemma:mixed-atoms}
Use the preceding excursion and pin assignments, and choose
$0<\kappa<c/2$, decreasing the positive exponent $c$ if necessary.
Then, for $X=Q,V$,
\[
 \max_P\left|X_P(f)-\mathcal A_X(P)\right|
                   \longrightarrow0
\]
in probability in the ordered limit.  The maximum is over the same
macroscopically reaching rims as above, on the preceding good events.
\end{lemma}

\begin{proof}
Let $\mathcal G=\mathcal G_{\eps,h}$ be the common event on which
the preceding qualified-pin assignments, density, winding bounds,
finite-cover conclusions, and simultaneous actual-arc comparisons
hold, together with bulk-patch agreement. Choose the finite-cover
budget tending to zero with $h$. The preceding results give
\[
 \limsup_{\eps\to0}\PP(\mathcal G^c)\le q(h),
                     \qquad q(h)\longrightarrow0.
\]
This event is distinct from the local usable-state events $G_s$.
Extend the adaptive coefficient assignment arbitrarily off $\mathcal G$;
all restricted counts and errors below are set to zero there.
We restrict the original probability measure by $\ind_{\mathcal G}$,
without dividing by $\PP(\mathcal G)$. Any additional mesh tests
at the finitely many levels used below have failure probability
tending to zero at fixed $h$ and fixed depth.

All discrepancies occur in a fixed compact neighborhood of
$\supp f$.  Tile it by $h$-squares and dyadically refine to
$s_j=2^{-j}h$.  At every level intersect the squares with the
deterministic partition recording membership in \emph{all candidate
cutoff balls}, with fixed disjoint Borel boundary conventions.  A
square meets only a bounded, menu-dependent number of these balls:
their radii are fixed multiples of $h$, and their centers lie on the
$h$-grid.  Thus there are at most a fixed number of atoms and children
per square.  The partitions are deterministic, regardless of which
pins were selected.

For a fixed rim call an occupied atom \emph{mixed} if two of its
tag centers have different assigned numerical coefficients.  Stop
unmixed atoms and refine mixed ones.  We first establish the
geometric implication that makes this refinement summable.  A mixed
atom has a nonzero coefficient involved, so its tags are near the
support.  Neither of two differently labeled tags can be in any
cutoff ball assigned to that rim.  Membership in every candidate
ball is constant throughout the atom, and an assigned ball near the
support assigns its entire rim mass a single coefficient.  Thus if
one tag belonged to it, both would have that coefficient.  This
argument concerns assigned balls; it asserts nothing about unused
candidates.

Choose a differently labeled tag with nonzero coefficient.  Since
it is outside all assigned cutoffs, it lies on a traversal chunk
between consecutive pins $p_i,p_{i+1}$.  The other tag cannot lie on
that same chunk.  In one direction along the rim between the two
visits the path must therefore pass $p_i$, and in the other direction
it must pass $p_{i+1}$.  This remains true if the second coefficient
is zero, or if its visit belongs to another excursion.  Each pin
core is within $h$ of its center, whereas the first tag is outside
both corresponding radius-$mh$ cutoffs.  Its distance from each core
is at least $(m_{\min}-1)h-o(h)$, with $m_{\min}>100$.  The two
visits in a common $s_j$-square are at distance $O(s_j)$.  They thus
have two distinct departures of order $h$, as illustrated in
Figure~\ref{fig:mixed-atom}.

\begin{figure}[tb]
\centering
\begin{tikzpicture}[scale=.75,>=Stealth,font=\small]
 \filldraw[fill=orange!12,draw=orange!75!black]
       (-.42,-.3) rectangle (.42,.3);
 \draw[very thick]
       (-.22,0)--(-1,1.2)--(-.85,2.5)--(.4,2.9)
       --(1.3,2.1)--(.9,1.0)--(.22,0)
       --(.9,-1.0)--(1.3,-2.1)--(.4,-2.9)
       --(-.85,-2.5)--(-1,-1.2)--cycle;
 \draw[dashed,blue!70!black] (-.85,2.5) circle (.48);
 \draw[dashed,blue!70!black] (-.85,-2.5) circle (.48);
 \fill (-.22,0) circle (2pt) (.22,0) circle (2pt);
 \fill (-.85,2.5) circle (2pt) (-.85,-2.5) circle (2pt);
 \node[above] at (-.22,.3) {$a$};
 \node[above] at (.22,.3) {$b$};
 \node[right,orange!60!black] at (.65,0) {$D$};
 \node[left] at (-1.5,2.5) {$p_i$};
 \node[left] at (-1.5,-2.5) {$p_{i+1}$};
 \draw[->,blue!70!black] (-1.6,.55)--(-1.9,1.45);
 \draw[->,blue!70!black] (-1.6,-.55)--(-1.9,-1.45);
 \node[left,align=right] at (-1.9,0)
       {both directions\\reach a pin core};
\end{tikzpicture}
\caption{The deterministic implication for a mixed atom $D$, shown in
orange.  The tag $a$
lies on the chunk between consecutive pins $p_i,p_{i+1}$; the other
tag cannot lie on that chunk.  Both intervening rim portions must
therefore reach a distant pin core.  Dashed circles depict the two
cutoffs.  The drawing is schematic and makes no regularity assertion
about a rim.}\label{fig:mixed-atom}
\end{figure}

For sufficiently fine fixed levels, Lemma~\ref{lemma:near-return}
gives the local necessary near-return arm event from scale $O(s_j)$
to a fixed multiple of $h$.  It can be tested inside, for example,
$B(y,10h)$ about the square center $y$, taking the outer radius of
the arm test below $8h$.  The same macroscopically reaching rim
supplies a one-arm event through a separated outer annulus, from
scale $32h$ to a fixed $\delta>0$, whose retained queries can be
confined outside $B(y,24h)$.  The numerical constants can be changed
to fit the fixed buffers.  At fixed $h,j$ the mesh exceptions vanish
and these tests use disjoint deterministic marked-Poisson
restrictions.  For each color $b$, the precise necessary-event
inclusion is
\[
 \begin{aligned}
 &\{\mathcal G;\ \text{atom mixed for some color-}b\text{ rim}\}
 \\
 &\qquad\subset E_{{\rm in},b}(s_j,h)\cap E_{{\rm out},b}(h,\delta)
                 \ \cup\ \{\text{mesh failure}\}.
 \end{aligned}
\]
The events on the right do not select a rim.  We multiply their
ordinary probabilities, not probabilities conditioned on adaptive
labels or on the good event.  The near-return gain and the one-arm
bound give, after decreasing $c>0$ to a common value,
\[
 \PP(E_{{\rm in},b}\cap E_{{\rm out},b})
                         \le C(s_j/h)^{2+c}h^c.
\]
The two colors cost only a factor two.  At the bounded number of
coarse levels where these numerical buffers do not fit, use the
one-arm estimate alone.

On $\mathcal G$, let $N_j$ count the geometric atoms to be considered at level $j$
because their parent is mixed for at least one relevant rim.  At
level zero count the occupied atoms in use.  Each atom is counted
once, regardless of how many rims visit it.  There are
$O_R(s_j^{-2})$ candidate squares, with a bounded number of atoms
and children each. Define $\widehat N_j=\ind_{\mathcal G}N_j$,
with value zero off $\mathcal G$. For $j\ge1$, apply the preceding
inclusion to the mixed parent at scale $s_{j-1}=2s_j$; its bounded
number of children and the factor two in scale change only the
constant. Summing the ordinary arm probabilities yields
\begin{equation}\label{actual:mixed-count}
 \limsup_{\eps\to0}\EE\widehat N_j
                 \le C_Rh^{-2+c}2^{-cj}.
\end{equation}
The same restricted bound at level zero follows from one arm.

Choose $f_D$ to be the value of $f$ at the center of the level-$j$
square containing $D$.  Define the nonnegative square sum
\begin{equation}\label{actual:bessel-level}
 B_j=\sum_{D,\lambda}\left(
       |X_\lambda(\ind_D)|^2+
       s_j^{-2}|X_\lambda(\ind_D(f-f_D))|^2\right).
\end{equation}
The sum includes every deterministic atom at this level and every
circuit-component label $\lambda$.  Since these are disjoint
deterministic masks, $|f-f_D|\le C_fs_j$, and all masks have support
in a fixed compact region, the Bessel bound
\eqref{fields:energy} and the local triangle moments give
\[
              \limsup_{\eps\to0}\EE\widehat B_j\le C_{f,R},
       \qquad \widehat B_j=\ind_{\mathcal H_\eps}B_j.
\]
Here $\mathcal H_\eps$ is the color-independent geometry event of
bulk-patch agreement, and is the whole space in the full-plane setting.
We have $\mathcal G\subset\mathcal H_\eps$. The full deterministic
atom and label sum defining $B_j$ is never restricted according to
the adaptive coefficients; its geometric indicator is allowed in
the conditional Bessel bound.
The bound is uniform in $j$ when that finite level is fixed before
the microscopic limit.

For one rim $P$ let $\mathcal S_j(P)$ be its stopped atoms at
level $j$, and write $\lambda(P)$ for its circuit-component label.
On such an atom its constant coefficient error relative to $f_D$
has magnitude at most $C_fr_h$.  The other error uses $f-f_D$.
Because $s_j\le h\le r_h$, Cauchy--Schwarz gives, on $\mathcal G$,
\begin{align*}
 &\left|\text{level-}j\text{ contribution to }
                    X_P(f)-\mathcal A_X(P)\right|\\
 &\quad\le C_fr_h\sqrt{\#\mathcal S_j(P)}
 \left[\sum_{D\in\mathcal S_j(P)}
       \bigl(|X_{\lambda(P)}(\ind_D)|^2+
       s_j^{-2}|X_{\lambda(P)}(\ind_D(f-f_D))|^2\bigr)\right]^{1/2}\\
 &\quad\le C_fr_h\sqrt{\widehat N_j\widehat B_j}.
\end{align*}
After extending the level error by zero off $\mathcal G$, the last
bound holds everywhere and bounds the maximum over all rims. In particular
no extra number-of-rims factor appears: $\widehat N_j$ counts geometric
atoms, while the second factor already sums over \emph{all} field
labels.  Taking expectation and applying Cauchy--Schwarz once more,
\eqref{actual:mixed-count} and \eqref{actual:bessel-level} bound the
limiting upper expectation by
\[
              C_{f,R}r_hh^{-1+c/2}2^{-cj/2}.
\]
This is summable over levels, with total
$O_{f,R}(r_hh^{-1+c/2})$.

To justify stopping, fix $h$ and a finite depth $J\ge1$. A mixed
level-$J$ atom has a mixed parent for the same rim, so it is itself
counted by $\widehat N_J$ on $\mathcal G$. Therefore
\[
 \limsup_{\eps\to0}
 \PP(\mathcal G\cap\{\text{some level-}J\text{ atom is mixed}\})
                  \le C_Rh^{-2+c}2^{-cJ}.
\]
On $\mathcal G$ when no terminal atom is mixed, the stopped atoms
through depth $J$ account for the whole discrepancy. The preceding
finite sum is bounded by its geometric series, independently of $J$.
For each $t>0$, Markov's inequality and the last display give
\[
 \begin{aligned}
 &\limsup_{\eps\to0}
   \PP\left(\max_P|X_P(f)-\mathcal A_X(P)|>t\right)\\
 &\qquad\le q(h)+C_Rh^{-2+c}2^{-cJ}
                +\frac{C_{f,R}}{t}r_hh^{-1+c/2}.
 \end{aligned}
\]
For a prescribed tolerance $\tau_h\downarrow0$, choose a finite
$J=J(h)$ making the middle term at most $\tau_h$, before taking
$\eps\to0$. Only finitely many mesh tests are needed at that fixed
$h,J$, so their failure terms vanish in this microscopic limit.
Finally let $h\to0$: both $q(h)+\tau_h$ and
$r_hh^{-1+c/2}=h^{c/2-\kappa}$ tend to zero. This proves the claim.
\end{proof}

\subsection{Completion of bulk transport}

Take $\kappa>0$ smaller than half the common positive exponent used
above.  Both terms in \eqref{actual:coefficient-cost} tend to zero,
since they equal, up to constants,
$h^{c-\kappa}(1+|\log h|^2)$ and $h^{c/2-\kappa}$.
Propositions~\ref{prop:dense-pins} and \ref{prop:real-arc}, followed
by Lemma~\ref{lemma:mixed-atoms}, show that for every $\xi>0$,
\begin{equation}\label{actual:ordered-transport}
 \lim_{h\to0}\limsup_{\eps\to0}
 \PP\left(\max_P|Q_P(f)-a_h^QV_P(f)|>\xi\right)=0.
\end{equation}
All limits are on the common microscopic subsequence chosen at the
start of the section.  The auxiliary finite covers and couplings
for a given patch and test require no further subsequence.

Pass along the $h$ sequence so that $a_h^Q\to a\in\C$.
The Bessel bound for $V$ implies
\[
 \limsup_{\eps\to0}
 \PP\left(\max_P|(a_h^Q-a)V_P(f)|>\xi\right)
                         \le C_f|a_h^Q-a|^2/\xi^2.
\]
Combine this with \eqref{actual:ordered-transport}, taking the
microscopic limsup at fixed $h$ and then $h\to0$.  This proves
\eqref{transport:bulk} for the resulting bounded scalar $a$, for
every fixed bulk patch and test $f$.  The scalar choices used only
the countable menu representatives and mean limits, and do not
depend on the later test or on points of evaluation of an observable.
The remaining assertion $\Rea a<1/2$ in
Theorem~\ref{thm:bulk-transport} is established in
Proposition~\ref{prop:nondegeneracy}.
